% Self-contained source. Compile twice with XeLaTeX into tmp/.
% Exercise wording is translated from the Chinese textbook, not an official English edition.
\documentclass[a4paper,12pt]{article}
\usepackage[margin=22mm,headheight=16pt,headsep=8mm,footskip=12mm]{geometry}
\usepackage{amsmath,amssymb,booktabs,enumitem,fancyhdr,xcolor,lastpage}
\usepackage[unicode,hidelinks]{hyperref}
\hypersetup{pdftitle={Stochastic Signal Analysis - Homework 2},pdfauthor={}}
\definecolor{muted}{gray}{0.38}
\pagestyle{fancy}\fancyhf{}
\fancyhead[L]{\small Stochastic Signal Analysis}
\fancyhead[R]{\small Homework 2}
\fancyfoot[L]{\scriptsize\color{muted}English translation from the Chinese textbook (5th ed.)}
\fancyfoot[R]{\small Page \thepage\ of \pageref*{LastPage}}
\renewcommand{\headrulewidth}{0.4pt}
\setlength{\parindent}{0pt}\setlength{\parskip}{0.5em}
\linespread{1.18}
\setlist[enumerate]{label=(\arabic*),leftmargin=2.3em,labelsep=0.4em,itemsep=0.35em,topsep=0.4em}
\newcommand{\problem}[3]{%
 {\small\color{muted}Chapter 2\enspace /\enspace Random Signals}\par
 {\small Issued: 13 October 2026\hfill Due: \textbf{27 October 2026}}\par
 \vspace{2mm}
 {\small Name: \underline{\hspace{35mm}}\hfill Student ID: \underline{\hspace{38mm}}}\par
 \vspace{4mm}
 {\LARGE\bfseries Exercise #1}\hfill{\normalsize\bfseries #2}\par
 {\small\color{muted}Textbook p.\ #3}\par\vspace{3mm}}
\newcommand{\answer}{\par\vspace{6mm}{\small\color{muted}Answer / Working}\par\vspace{1mm}{\color{muted}\hrule height 0.3pt}\vfill}
\begin{document}
\problem{2.3}{Required}{61}

A coin is tossed, and a random process is defined by
\[
X(t)=\begin{cases}
\cos\pi t,&\text{if heads occurs},\\
2t,&\text{if tails occurs}.
\end{cases}
\]
Heads and tails are equally likely. Find:
\begin{enumerate}
\item The one-dimensional distribution functions $F_X(x,1/2)$ and $F_X(x,1)$ of $X(t)$;
\item The two-dimensional distribution function $F_X(x_1,x_2;1/2,1)$ of $X(t)$;
\item Sketch the graphs of the above distribution functions.
\end{enumerate}
\answer

\newpage
\problem{2.13}{Required}{62}

Consider the sinusoidal voltage signal
\[
X(t)=A\cos(\omega t+\Theta),
\]
where $\omega$ is a constant, $A\sim U(-1,+1)$, and $\Theta\sim U(-\pi,+\pi)$. The random variables $A$ and $\Theta$ are independent.

If this voltage is applied across a parallel RC circuit, find the total current signal and its mean-square value.
\answer

\newpage
\problem{2.16}{Required}{62}

Let the random process be
\[
X(t)=A\cos\omega t+B\sin\omega t,
\]
where $\omega$ is a constant and $A$ and $B$ are independent Gaussian random variables. Given that
\[
E[A]=E[B]=0,\qquad E[A^2]=E[B^2]=\sigma^2,
\]
find the one-dimensional and two-dimensional probability density functions of $X(t)$.
\answer

\newpage
\problem{2.19}{Required}{62}

Let the random vector $(X,Y)\sim N(\boldsymbol{\mu},\boldsymbol{C})$, where
\[
\boldsymbol{\mu}=\begin{bmatrix}2\\2\end{bmatrix},\qquad
\boldsymbol{C}=\begin{bmatrix}2&3\\3&5\end{bmatrix}.
\]
Find the joint probability density function of $(X,Y)$ and its characteristic function $\phi_{XY}(u,v)$.
\answer

\newpage
\problem{2.1}{Optional}{61}

Consider the sinusoidal random signal $X(t)=V\cos\omega t$, where $0\le t<\infty$, $\omega$ is a constant, and $V$ is uniformly distributed on $[0,1]$.
\begin{enumerate}
\item Sketch two sample functions of this process.
\item Determine the probability density function of the random variable $X(t_i)$ at each of the times
\[
t_i=0,\quad\frac{\pi}{4\omega},\quad\frac{3\pi}{4\omega},\quad\frac{\pi}{\omega}.
\]
\item Find the probability density function of $X(t_0)$ when $t_0=\pi/(2\omega)$.
\end{enumerate}
\answer

\newpage
\problem{2.4}{Optional}{61}

Suppose that the binary data sequence $\{B(n),n=1,2,3,\ldots\}$ is a Bernoulli random sequence. Each bit corresponds to a random variable $B(n)$, with
\[
P[B(n)=0]=0.2,\qquad P[B(n)=1]=0.8.
\]
\begin{enumerate}
\item What is the probability that four consecutive bits form the word $\{1011\}$?
\item What is the mean word formed by four consecutive bits?
\item Which word formed by four consecutive bits has the highest probability?
\end{enumerate}
\answer

\newpage
\problem{2.7}{Optional}{61--62}

The position of a moving particle is described by the straight-line process
\[
X(t)=Vt+X_0,
\]
where $V\sim N(1,1)$ and $X_0\sim N(0,2)$ are independent.

Find the one-dimensional probability density function, mean, and variance of the random variable $X(t)$ at time $t$.
\answer

\newpage
\problem{2.8}{Optional}{62}

Suppose that the Bernoulli sequence $\{I_n,n=1,2,\ldots\}$ takes the values $-1$ and $+1$ with equal probabilities.

Find its one-dimensional probability density function, mean, and covariance function.
\answer

\newpage
\problem{2.11}{Optional}{62}

Consider the two random signals
\[
X(t)=A\sin(\omega t+\Theta),\qquad
Y(t)=B\cos(\omega t+\Theta),
\]
where $A$ and $B$ are random variables with unspecified distributions, $\Theta$ is uniformly distributed on $[0,2\pi]$, and $\omega$ is a constant. The random variables $A$, $B$, and $\Theta$ are pairwise statistically independent.
% Preserve the textbook's pairwise independence assumption; do not strengthen it to mutual independence.
\begin{enumerate}
\item Find the cross-correlation function $R_{XY}(t_1,t_2)$ of the two random signals.
\item Discuss their orthogonality, uncorrelatedness, and statistical independence.
\end{enumerate}
\answer

\newpage
\problem{2.18}{Optional}{62}

A Gaussian signal has mean and covariance function
\[
m_X(t)=2,\qquad C_X(t_1,t_2)=8\cos(t_1-t_2).
\]
Write down its three-dimensional probability density function at the times
\[
t_1=0,\qquad t_2=0.5,\qquad t_3=1.
\]
% Preserve the textbook question. A later solution must address the singular covariance matrix.
\answer
\end{document}
